% C.3ar: restricted positive theorem; not a closure of the unrestricted top gate.
\paragraph{C.3ar (Proved in the stated product-of-curves scope): actual families kill the absolute top terms.}
Let $d\ge2$. For $0\le i<d$, let $C_i/\overline{\mathbb Q}$ be
smooth projective geometrically connected curves and let $L_i$ be
very ample line bundles over $\overline{\mathbb Q}$. Put
\[
 V=\left(\prod_{i=0}^{d-1}C_i\right)_{\mathbb C},
 \qquad L=\boxtimes_{i=0}^{d-1}L_i,
 \qquad S=\bigoplus_{n\ge0}H^0(V,L^n).
\]
Assume $S$ is normal and generated in degree one. Let $R$ be the
strict henselization of $S_{S_+}$ at its complex vertex, let
$X=\operatorname{Spec}R\setminus\{\mathfrak m_R\}$, and let
$Z\to\operatorname{Spec}R$ be its regular cone resolution, with
exceptional divisor $i:V\hookrightarrow Z$ and projection $p:Z\to V$.
Then, for ordinary algebraic equivalence on actual fixed products,
\begin{equation}\label{eq:C3ar-answer}
 B^d(Z)_{\mathbb Q}\xrightarrow[\deg\circ i^!]{\ \sim\ }\mathbb Q,
 \qquad B^d(X)_{\mathbb Q}=0.
\end{equation}
The coefficient field here is $\mathbb C$, not
$\overline{\mathbb Q}$. The conclusion concerns rational ordinary
quotients; it neither asserts integral finiteness nor settles arbitrary
complex cones or arbitrary isolated local singularities.

\begin{proof}
\emph{1. Finite detection and the exact relative group.}
Write $V_a=Z\times_R R/\mathfrak m_R^a$. As in C.3ab--ad,
\[
 V_a=\operatorname{Spec}_V
       (\mathcal O_V\oplus L\oplus\cdots\oplus L^{a-1}),
 \qquad K_0(Z)\longrightarrow K_0(V_a)
                   \text{ is injective for }a\gg0.
\]
The rational Chern character identifies
$\mathrm{CH}^d(Z)_{\mathbb Q}$ with $K_0^{(d)}(Z)_{\mathbb Q}$;
the finite pointed-\'etale models used in C.3ad justify this assertion
and all its pullback compatibilities for this $Z$.

The relative Adams calculation of C.3ad is valid over $\mathbb C$
before its final vanishing step. It gives
\begin{equation}\label{eq:C3ar-relative}
 K_0^{(d)}(V_a,V)_{\mathbb Q}
   \simeq\bigoplus_{j=1}^{a-1}
       H^d(V,\Omega^{d-1}_{V/\mathbb Q}\otimes L^j).
\end{equation}
We recall why this is an isomorphism, not merely a possible graded
subquotient. The relative sheaves of weight $d$ vanish in degrees
$q<d$, while the degree-$d$ sheaf is
$\bigoplus_{1\le j<a}\Omega^{d-1}_{V/\mathbb Q}\otimes L^j$.
On the diagonal of the descent spectral sequence, only $(d,d)$
can survive. An incoming differential of length $r\ge2$ would
start in sheaf degree $d-r+1<d$; an outgoing one would have
cohomological degree greater than $d$. Both vanish. The construction
is natural for restriction in $a$ and respects the internal degree
$j$ of the split cone. The exact inputs are the arbitrary
commutative-$\mathbb Q$-algebra calculation of
\citet[Theorems 4.6.7--4.6.8, pp.~151--152]{D_Loday1992} and
the scheme-compatible relative Chern character of
\citet[Theorem 6.1, Corollary 6.2 and Theorem B.10,
author pp.~15 and 25]{D_CortinasHaesemeyerWeibel2009}.

Because $V$ has the displayed model over $\overline{\mathbb Q}$,
its absolute differential sequence has the canonical splitting
\[
 \Omega^1_{V/\mathbb Q}
  =\Omega^1_{V/\mathbb C}\oplus
                 (\mathcal O_V\otimes_{\mathbb C}
                           \Omega^1_{\mathbb C/\mathbb Q}).
\]
Kodaira--Akizuki--Nakano vanishing kills
$H^d(V,\Omega^q_{V/\mathbb C}\otimes L^j)$ for $q>0$, $j>0$.
Consequently \eqref{eq:C3ar-relative} becomes
\begin{equation}\label{eq:C3ar-absolute}
 K_0^{(d)}(V_a,V)_{\mathbb Q}
 \simeq\bigoplus_{j=1}^{a-1}
       H^d(V,L^j)\otimes_{\mathbb C}
                         \Omega^{d-1}_{\mathbb C/\mathbb Q}.
\end{equation}
These are identities of underlying rational vector spaces, with their
indicated complex-linear descriptions. We do not infer that a Chow
group is a complex vector space. A finite affine cover justifies
commuting cohomology with the possibly infinite-dimensional constant
space of differentials.

\emph{2. Every chosen Picard deformation is an actual line bundle.}
For $j\ge1$ and $a\in H^1(V,L^j)$, consider the formal line bundle
whose logarithmic class is $a$ in homogeneous degree $j$ and zero
in every other positive degree. More concretely, on a trivializing
cover its transition functions are $\exp(t^j a_{\alpha\beta})$,
where $t$ denotes a local conormal coordinate. The transformation
rule for $t$ is exactly the one needed to glue the $L^j$-valued
\v Cech cocycle. Formal logarithm identifies the kernel of
restriction to $V$ with
\[
       \prod_{n\ge1}H^1(V,L^n)
           =\bigoplus_{n\ge1}H^1(V,L^n),
\]
where the second equality follows from Serre vanishing.

We use the actual line-bundle argument of C.3ao, Step~2, for this
particular cone. Formal existence on the proper completed resolution
gives a line bundle there. Restricting to the completed puncture and
using the exact Picard comparison of Elkik gives a line bundle on
$X$. It extends as a Cartier divisor to the regular scheme $Z$.
Its completion differs from the prescribed formal line bundle by an
integral multiple of the exceptional divisor, because that divisor
is the entire reduced complement of the puncture. Twisting by that
multiple removes the difference. We obtain an actual line bundle
$A_{j,a}$ on $Z$ such that
\begin{equation}\label{eq:C3ar-line}
 A_{j,a}|_V\simeq\mathcal O_V,
 \qquad \log(A_{j,a}|_{\widehat Z})=t^j a.
\end{equation}
All hypotheses in this argument are those checked in C.3ao: the
cone is normal, its puncture is regular, and its regular resolution
has the single smooth exceptional Cartier divisor $V$. This
constructs actual Picard classes; it does not assert algebraization
of arbitrary formal $K_0$ classes.

\emph{3. The leading product and its normalization.}
For line bundles $\lambda_i\in\operatorname{Pic}^0(C_i)(\mathbb C)$,
$1\le i<d$, put
\begin{equation}\label{eq:C3ar-product}
 \beta(j,a;\lambda_1,\ldots,\lambda_{d-1})
  =c_1(A_{j,a})\prod_{i=1}^{d-1}p_i^*c_1(\lambda_i)
                      \in\mathrm{CH}^d(Z),
\end{equation}
where $p_i$ is the composite $Z\to V\to(C_i)_{\mathbb C}$.
Its restriction to $V$ is zero. If each $\lambda_i$ is a point
difference $\mathcal O_{C_i}(q_i-o_i)$, this cycle is ordinarily
algebraically trivial: vary $q_1$ on the actual smooth projective
curve $(C_1)_{\mathbb C}$, keep the other factors fixed, and use the
diagonal-minus-constant divisor. Finite sums of these products have
the same property.

Let
\[
 \epsilon_i(\lambda_i)\in
 H^1(C_i,\mathcal O_{C_i})\otimes_{\mathbb C}
                        \Omega^1_{\mathbb C/\mathbb Q}
\]
be the absolute first Chern class. The relative first Chern class
in $H^1(C_i,\Omega^1_{C_i/\mathbb C})$ is zero because
$\deg\lambda_i=0$, so there is no omitted relative component.
Under \eqref{eq:C3ar-absolute}, the first possible internal degree
of \eqref{eq:C3ar-product} is $j$, and its degree-$j$ term is,
with the usual ordered cup-product convention,
\begin{equation}\label{eq:C3ar-leading}
       a\smile\epsilon_1(\lambda_1)\smile\cdots
                   \smile\epsilon_{d-1}(\lambda_{d-1}).
\end{equation}
An overall Koszul sign is immaterial and can be absorbed into $a$.
There is no zero or unspecified rational multiplier in this formula.

Here is a direct normalization check. On a trivializing affine
cover, the leading transition for $A_{j,a}$ is $1+t^j a$.
The product of the first Chern cocycles is the Milnor symbol
\[
                  \{1+t^j a,g_1,\ldots,g_{d-1}\}.
\]
Its logarithmic cyclic representative is
$t^j a\,d\log g_1\wedge\cdots\wedge d\log g_{d-1}$.
Applying the Connes differential gives exactly the $d\log$ of the
displayed symbol. The positive-degree horizontal representative is
unique by the calculation in C.3ad; equivalently its differential
is injective in positive internal degree, by contraction with the
Euler derivation and division by $j$. This proves the coefficient
in \eqref{eq:C3ar-leading}. The exact symbol formula is
\citet[Lemma 2.2 and Proposition 2.3, author pp.~3--4]{D_BlochEsnaultKerz2014};
their Proposition 6.1 proof, author pp.~12--13, explicitly checks
that the canonical comparison followed by the differential is
$d\log$. Multiplicativity of the scheme-compatible Chern character
is also explicit in
\citet[author pp.~16 and 25]{D_CortinasHaesemeyerWeibel2009}.
Thus the same calculation applies to the cup products in descent.
We use this leading-term formula at each finite stage. We do not
replace the pro-isomorphisms in Bloch--Esnault--Kerz, Section~6,
by unasserted isomorphisms at a fixed thickening. Equivalently, a
fixed leading term may be checked after restriction from an arbitrarily
higher compatible stage of these actual Picard products; the
fixed-level weight-$d$ calculation itself is supplied separately
by Loday and Corti\~nas--Haesemeyer--Weibel.

\emph{4. Fibre evaluations span the coefficient forms, with actual coefficients.}
Fix $i\ge1$, and choose a nonconstant morphism
$u_i:C_i\to\mathbb P^1$ over $\overline{\mathbb Q}$ of degree
$e_i>2g(C_i)-2$. Such a morphism exists by Riemann--Roch, after
choosing a sufficiently large pole divisor over the algebraically
closed field $\overline{\mathbb Q}$. Let $b\in\mathbb C$ be
transcendental over $\overline{\mathbb Q}$. The fibre $u_i^{-1}(b)$
is reduced, since every branch value is algebraic. Choose an
algebraic base point $o_i\in C_i(\overline{\mathbb Q})$.
For each $q\in u_i^{-1}(b)$, the residue calculation of C.3ao gives
\begin{equation}\label{eq:C3ar-evaluation}
 \left\langle\epsilon_i\bigl(\mathcal O(q-o_i)\bigr),\omega\right\rangle
   =-q^*\omega
   =-\left(\frac{\omega}{du_i}\right)(q)\,db
       \quad\bigl(\omega\in H^0(C_i,K_{C_i})\bigr).
\end{equation}
The model over $\overline{\mathbb Q}$ is essential for this formula:
all differentials here are absolute over $\mathbb Q$ and the
base point contributes zero.

The functionals obtained by evaluation at the fibre points span
the full dual of $H^0(C_i,K_{C_i})$. Otherwise a nonzero holomorphic
form would vanish at every one of the $e_i$ distinct unramified
points. This contradicts the canonical degree:
\[
                    \deg K_{C_i}=2g(C_i)-2<e_i.
\]
By Serre duality the complex-linear span of the actual classes
in \eqref{eq:C3ar-evaluation} is therefore
\begin{equation}\label{eq:C3ar-fibre-span}
                  H^1(C_i,\mathcal O_{C_i})\otimes db.
\end{equation}
This is a finite-dimensional linear-algebra statement about
coefficient forms, not an assertion that the ordinary Chow subgroup
is closed under multiplication by complex numbers.

For every $j>0$ choose nonzero sections
$\sigma_i\in H^0(C_i,L_i^j)$, $i\ge1$. Cup product with $\sigma_i$
is surjective:
\[
 H^1(C_i,\mathcal O_{C_i})\xrightarrow{\,\sigma_i\,}
                           H^1(C_i,L_i^j).
\]
Its Serre-dual map is multiplication by $\sigma_i$ from
$H^0(C_i,K_{C_i}\otimes L_i^{-j})$ to $H^0(C_i,K_{C_i})$,
which is injective. K\"unneth gives
\[
             H^d(V,L^j)=\bigotimes_{i=0}^{d-1}H^1(C_i,L_i^j).
\]
Every absolute $(d-1)$-form is a finite sum of
$c\,db_1\wedge\cdots\wedge db_{d-1}$; terms with an algebraic
$b_i$ vanish. To obtain a specified decomposable coefficient
$v_0\otimes\cdots\otimes v_{d-1}$ times such a form, lift
$v_i$ through the displayed cup maps for $i\ge1$, and express
these lifts using \eqref{eq:C3ar-fibre-span}. In each resulting
finite product take
\[
         a=c'\,v_0\otimes\sigma_1\otimes\cdots
                                      \otimes\sigma_{d-1}
                              \in H^1(V,L^j),
\]
where $c'$ is the product of the complex coefficients of that
term, including $c$. Construct the actual line bundle $A_{j,a}$
with this $a$. Formula \eqref{eq:C3ar-leading} then realizes the
specified term by a finite sum of actual ordinary cycles
\eqref{eq:C3ar-product}. No complex scalar is applied to a Chow
class. If one of the relevant $H^1(C_i,L_i^j)$ is zero, the
coefficient space itself is zero and no construction is needed.

\emph{5. Finite successive elimination and the ordinary quotient.}
Fix a detecting $a_0$ from Step~1. Let
$z\in\mathrm{CH}^d(Z)_{\mathbb Q}$ satisfy $i^!z=0$.
Its image in $K_0^{(d)}(V_{a_0})_{\mathbb Q}$ lies in the
relative group \eqref{eq:C3ar-absolute}. At internal degree one,
Step~4 supplies a finite sum of actual ordinary cycles with the
same coefficient. Subtract it. Continue through degrees
$2,\ldots,a_0-1$. A product constructed from $A_{j,a}$ has no
term in a smaller degree, so earlier cancellations are preserved.
There are only finitely many steps. The remaining class has zero
image in $K_0(V_{a_0})_{\mathbb Q}$ and hence is zero in
$K_0(Z)_{\mathbb Q}$ by finite detection. The rational Chern
character on $Z$ shows that it is zero in
$\mathrm{CH}^d(Z)_{\mathbb Q}$ as well. We have proved
\[
             \ker(i^!:\mathrm{CH}^d(Z)_{\mathbb Q}
                            \to\mathrm{CH}_0(V)_{\mathbb Q})
                      \subseteq\mathrm{CH}^d(Z)_{\mathrm{alg},\mathbb Q}.
\]
For arbitrary $z$, subtract $p^*i^!z$ first. Since $p^*$ and
$i^!$ preserve ordinary equivalence on actual products and
$i^!p^*=1$, the ordinary quotient on $Z$ is the ordinary
zero-cycle quotient on the connected projective variety $V$.
It is $\mathbb Q$, via degree. This proves the first assertion
of \eqref{eq:C3ar-answer}.

Finally, Chow localization gives
$\mathrm{CH}^{d-1}(V)\to\mathrm{CH}^d(Z)\to\mathrm{CH}^d(X)\to0$.
Self-intersection identifies the first map, after restriction to
$V$, with multiplication by $-c_1(L)$. Some integral curve in
$V$ has positive $L$-degree. Its image therefore generates the
one-dimensional rational ordinary quotient of $Z$.
The quotient on $X$ is zero, as claimed. Every ordinary cycle
used in this proof came from a fixed proper curve parameter;
no family on a henselized product or an arbitrary formal
$K_0$ object was substituted for an actual family.
\end{proof}
